
\medskip
\phantomsection\label{par:sga3-VIA-5.2}%
\noindent\textbf{5.2.---}
When the homomorphism of \(A\)-groups \(u:F\to G\) is a
monomorphism, the result of \hyperref[par:sga3-VIA-5.1]{5.1} may also be recovered from the
results of Expos\'e~IV. Indeed, the canonical morphism
\[
  p:G\longrightarrow F\backslash G
\]
is faithfully flat and open by \hyperref[thm:sga3-VIA-3.2]{3.2}; it is therefore a covering for
the \((\mathrm{fpqc})\) topology by \SGARefLink{sga3:iv:proposition:6:3:1}{Expos\'e~IV, 6.3.1}, and we may
apply \SGARefLink{sga3:iv:corollary:5:2:2}{Expos\'e~IV, Corollaries~5.2.2} and 5.2.4.

In particular, suppose, in addition to the hypotheses of \hyperref[thm:sga3-VIA-3.2]{3.2}, that
\(u\) is the inclusion in \(G\) of a normal subgroup \(F\). Then
there is a unique \(A\)-group structure on \(F\backslash G\) for
which the canonical morphism
\[
  p:G\longrightarrow F\backslash G
\]
is a homomorphism of \(A\)-groups.%
\footnote{\label{note:sga3-VIA-51}Editor's note: we have added the
following sentence.}
This proves \hyperref[concl:sga3-VIA-3.3.2-v]{Theorem~3.3.2(v)}.
